\documentclass[aps,prb,onecolumn]{revtex4-2}

\usepackage{amsmath,amssymb,amsfonts}
\usepackage{graphicx}
\usepackage{hyperref}
\usepackage{bm}
\usepackage{booktabs}
\usepackage{longtable}

\newcommand{\starG}{\star_G}
\newcommand{\Ghat}{\widehat{G}}
\newcommand{\CC}{\mathbb{C}}
\newcommand{\RR}{\mathbb{R}}
\newcommand{\ZZ}{\mathbb{Z}}

\begin{document}

\title{Supplementary Material:\\
  Full Spacegroup-Fourier Spectroscopy for\\
  Superconductivity Prediction via the $\starG$ Transform}

\author{Lior Horesh}
\affiliation{MIT}

\date{\today}

\maketitle

%% ====================================================================
\section{Character tables for the 11 proper point groups}
\label{sec:irrep-tables}

The 11 proper (chiral) crystallographic point groups and their
irreducible representations are listed below. Dimension $d$,
angular momentum quantum number $\ell$, and Frobenius-Schur
indicator $\nu$ are given for each irrep.

\begin{table}[h]
\caption{Irrep summary for proper crystallographic point groups.}
\begin{ruledtabular}
\begin{tabular}{llllll}
Group & $|G|$ & Irreps & $d$ & $\ell$ & $\nu$ \\
\midrule
C$_1$ & 1  & A       & 1       & 0       & $+1$ \\
\midrule
C$_2$ & 2  & A       & 1       & 0       & $+1$ \\
      &    & B       & 1       & 1       & $+1$ \\
\midrule
C$_3$ & 3  & A       & 1       & 0       & $+1$ \\
      &    & E       & 2       & 1       & $0$  \\
\midrule
C$_4$ & 4  & A       & 1       & 0       & $+1$ \\
      &    & B       & 1       & 2       & $+1$ \\
      &    & E       & 2       & 1       & $0$  \\
\midrule
C$_6$ & 6  & A       & 1       & 0       & $+1$ \\
      &    & B       & 1       & 3       & $+1$ \\
      &    & E$_1$   & 2       & 1       & $0$  \\
      &    & E$_2$   & 2       & 2       & $0$  \\
\midrule
D$_2$ & 4  & A       & 1       & 0       & $+1$ \\
      &    & B$_1$   & 1       & 1       & $+1$ \\
      &    & B$_2$   & 1       & 1       & $+1$ \\
      &    & B$_3$   & 1       & 1       & $+1$ \\
\midrule
D$_3$ & 6  & A$_1$   & 1       & 0       & $+1$ \\
      &    & A$_2$   & 1       & 1       & $+1$ \\
      &    & E       & 2       & 1       & $+1$ \\
\midrule
D$_4$ & 8  & A$_1$   & 1       & 0       & $+1$ \\
      &    & A$_2$   & 1       & 1       & $+1$ \\
      &    & B$_1$   & 1       & 2       & $+1$ \\
      &    & B$_2$   & 1       & 2       & $+1$ \\
      &    & E       & 2       & 1       & $+1$ \\
\midrule
D$_6$ & 12 & A$_1$   & 1       & 0       & $+1$ \\
      &    & A$_2$   & 1       & 1       & $+1$ \\
      &    & B$_1$   & 1       & 3       & $+1$ \\
      &    & B$_2$   & 1       & 3       & $+1$ \\
      &    & E$_1$   & 2       & 1       & $+1$ \\
      &    & E$_2$   & 2       & 2       & $+1$ \\
\midrule
T     & 12 & A       & 1       & 0       & $+1$ \\
      &    & E       & 2       & 1       & $0$  \\
      &    & T       & 3       & 1       & $+1$ \\
\midrule
O     & 24 & A$_1$   & 1       & 0       & $+1$ \\
      &    & A$_2$   & 1       & 1       & $+1$ \\
      &    & E       & 2       & 2       & $+1$ \\
      &    & T$_1$   & 3       & 1       & $+1$ \\
      &    & T$_2$   & 3       & 1       & $+1$ \\
\end{tabular}
\end{ruledtabular}
\end{table}

%% ====================================================================
\section{Spacegroup to point group mapping}
\label{sec:spg-mapping}

Each of the 230 crystallographic spacegroups maps to one of 11
proper point groups. The distribution is:

\begin{table}[h]
\caption{Spacegroup count per proper point group.}
\begin{ruledtabular}
\begin{tabular}{lrl}
Point group & Count & Example spacegroups \\
\midrule
D$_2$ & 49 & 16--24 (orthorhombic) \\
C$_2$ & 31 & 3--5 (monoclinic) \\
D$_4$ & 30 & 89--98 (tetragonal) \\
C$_4$ & 26 & 75--80 (tetragonal) \\
T     & 18 & 195--199 (cubic) \\
O     & 18 & 207--214 (cubic) \\
D$_3$ & 17 & 149--155 (trigonal) \\
C$_3$ & 13 & 143--146 (trigonal) \\
C$_6$ & 12 & 168--173 (hexagonal) \\
D$_6$ & 10 & 177--182 (hexagonal) \\
C$_1$ &  6 & 1 (triclinic) \\
\end{tabular}
\end{ruledtabular}
\end{table}

%% ====================================================================
\section{Cross-validation: full-spectrum $\chi_3$ vs.\ v14 z3}
\label{sec:cross-validation}

We compared our fingerprint-derived $\chi_3$ (sum of non-trivial
spectral weights from the full angular spectrum observable) against
the v14 predictor's z3 (2D angular projection restricted to
$\ZZ/3$) for 735 matched materials.

\begin{table}[h]
\caption{Per-point-group cross-validation statistics.}
\begin{ruledtabular}
\begin{tabular}{lrrrr}
PG & $n$ & Mean diff. & Std diff. & Pearson $r$ \\
\midrule
C$_1$ &  11 & $-0.148$ & 0.172 & N/A \\
C$_2$ &  26 & $-0.132$ & 0.221 & 0.245 \\
C$_3$ &   3 & $-0.021$ & 0.151 & 0.930 \\
C$_4$ &  62 & $-0.100$ & 0.287 & $-0.577$ \\
C$_6$ &   1 & $-0.024$ & 0.000 & N/A \\
D$_2$ &  59 & $-0.134$ & 0.275 & $-0.151$ \\
D$_3$ &  50 & $+0.063$ & 0.268 & 0.385 \\
D$_4$ &  78 & $-0.068$ & 0.219 & 0.260 \\
D$_6$ &  99 & $+0.055$ & 0.293 & 0.196 \\
T     &  33 & $+0.301$ & 0.193 & 0.915 \\
O     & 313 & $+0.241$ & 0.216 & 0.943 \\
\midrule
All   & 735 & $+0.094$ & 0.287 & 0.626 \\
\end{tabular}
\end{ruledtabular}
\end{table}

The moderate overall correlation ($r = 0.63$) is expected because the
two quantities use different observables (3D bond geometry vs.\ 2D
angular projection) and different group restrictions (full point group
vs.\ $\ZZ/3$). The strong correlation within O and T groups reflects
the fact that for these cubic groups, the $\ZZ/3$ cyclic subgroup
captures a substantial portion of the full group structure.

%% ====================================================================
\section{Frobenius-Schur indicators}
\label{sec:fs-indicators}

The Frobenius-Schur indicator $\nu(\rho) = (1/|G|)\sum_g \chi(g^2)$
classifies each irrep:

\begin{itemize}
  \item $\nu = +1$: real type (orthogonal). All 1D irreps of
    crystallographic point groups, and all E-type irreps of D$_n$
    and O groups.
  \item $\nu = 0$: complex type (unitary, not self-conjugate). The
    E-type irreps of C$_3$, C$_4$, C$_6$, and T.
  \item $\nu = -1$: quaternionic type (symplectic). Does not occur
    among the 11 proper crystallographic point groups.
\end{itemize}

For the Plancherel identity with real matrix representations, the
Frobenius-Schur correction multiplies the spectral contribution of
complex-type irreps by $1/2$. This correction is necessary because
the realification of a conjugate pair of complex 1D irreps produces
a 2D real matrix whose Frobenius norm counts each degree of freedom
twice.

Groups requiring the correction: C$_3$ (E irrep), C$_4$ (E irrep),
C$_6$ (E$_1$ and E$_2$ irreps), and T (E irrep). The correction does
not affect the normalized spectral weights or $\chi_3$ values, since
the factor cancels in the ratio~(Eq.~3 of the main text).

%% ====================================================================
\section{Lean 4 proof listing}
\label{sec:lean-proofs}

The formal verification comprises four Lean~4 modules:

\begin{enumerate}
  \item \texttt{Plancherel.lean}: Plancherel identity from column
    orthogonality via the kernel trick. Defines \texttt{UnitaryRep},
    \texttt{IrrepFamily}, \texttt{frobSq}, \texttt{fourierCoeff},
    and proves \texttt{spectral\_weights\_sum} by expanding the
    Frobenius norm as a double sum, rearranging to put the irrep
    index innermost, applying column orthogonality, and collapsing
    to the $L^2$ norm.

  \item \texttt{SpectralWeights.lean}: Normalization theorems.
    Proves \texttt{normalized\_sum\_eq\_one} (weights sum to 1),
    \texttt{normalized\_nonneg} (each weight $\ge 0$),
    \texttt{normalized\_le\_one} (each weight $\le 1$), and
    the combined \texttt{spectral\_fingerprint\_is\_probability}.

  \item \texttt{FrobeniusSchur.lean}: FS indicator definition and
    correction factors. Proves \texttt{correction\_real\_type}
    ($c = 1$ for $\nu = 1$) and \texttt{correction\_complex\_type}
    ($c = 1/2$ for $\nu = 0$). The boundedness theorem
    (\texttt{fsIndicator\_bounded}) is stated but left for
    future formalization, as it requires Schur orthogonality
    in its full generality.

  \item \texttt{Chi3Reduction.lean}: $\chi_3$ as complement of
    trivial weight sum. Proves \texttt{chi3\_complement} and
    \texttt{chi3\_determined\_by\_fingerprint}. The strict
    refinement existence theorem is stated constructively.
\end{enumerate}

All proofs are available in the \texttt{lean/GSpacegroup/} directory
of the accompanying repository.

\end{document}
